\scanpage{92}\mspage{91}

La condition At 14 (peut-être \unsure{facultative}) semble \deleted{plus} un peu \unsure{abracadabrant|abracadabrante}, traduite en termes de \unsure{sous-atelier|sous-ateliers} d'un $\mathrm{Fig}(\mathscr L)$ -- ce n'est plus dans la nature d'une condition de stabilité.

De ce point de vue, l'axiome At 14 bis \unsure{``|``(}facultatif) est plus sympathique :

\uline{At ens fac} 5 : Soit $X\in M$, et $F'\preccurlyeq\partial X$

\noindent\begin{tabular}{@{}||p{.90\linewidth}@{}}
avec $ {\unsure{\mathbb F|F}}'\in M$. Considérons \unsure{l'unique|le 'unique} figure \deleted{\unsure{en}} $F\in\mathrm{Fig}(L)$ qui soit élémentaire et telle que $\partial F=F'$. Alors $F\in M$.
\end{tabular}

Mais ce sera pas vrai pour les ateliers simpliciaux !

\begin{center}\rule{1.7cm}{.35pt}\end{center}

Passons à la situation d'un atelier $\mathcal F$, et d'un sous-atelier $\mathcal F'\subset\mathcal F$.

(cf conditions Sous at 1, 2 p. \unsure{47|41})

Les conditions At 1 à At 7 sur \deleted{\unsure{tout|tous}} $\mathcal F$ se transportent aussitôt à $\mathcal F'$.

Pour les conditions impliquant $L$ (At 8, 9, $9_{\Sigma}$, 10) il faut déjà pouvoir comparer $\deleted{\mathcal F}\unsure{L|\mathscr L}$ et $\deleted{\mathcal F'}L'$. \unsure{Èvidemment|Évidemment} \deleted{$\unsure{L}$}
\[
L\cap\mathcal F'\subset L'
\]
